% note.tex
% SLE Hadamard State on Bianchi II / Heisenberg Nilmanifold
% Sequel to Banerjee-Niedermaier 2023 (arXiv:2305.11388)
% Draft 2026-05-03
%
% WARNING: Citation Pukanszky "TAMS 126, 487-507" is WRONG.
%          Correct cite: J. Funct. Anal. 1 (1967) 255-280. See gaps.md.
%
\documentclass[11pt, a4paper]{article}

\usepackage{amsmath, amssymb, amsthm}
\usepackage{mathtools}
\usepackage{hyperref}
\usepackage{enumitem}
\usepackage[margin=2.5cm]{geometry}

% Theorems
\newtheorem{theorem}{Theorem}
\newtheorem{proposition}[theorem]{Proposition}
\newtheorem{lemma}[theorem]{Lemma}
\newtheorem{corollary}[theorem]{Corollary}
\newtheorem{definition}[theorem]{Definition}
\newtheorem{remark}[theorem]{Remark}

% Notation
\newcommand{\RR}{\mathbb{R}}
\newcommand{\ZZ}{\mathbb{Z}}
\newcommand{\CC}{\mathbb{C}}
\newcommand{\NN}{\mathbb{N}}
\newcommand{\HH}{\mathfrak{h}}  % Heisenberg Lie algebra
\newcommand{\Hgroup}{H_3}        % Heisenberg group
\newcommand{\Nil}{\mathrm{Nil}^3}
\newcommand{\piom}{\pi_\omega}
\newcommand{\dpiom}{d\pi_\omega}
\newcommand{\WF}{\mathrm{WF}}
\newcommand{\HS}{\mathrm{HS}}
\newcommand{\Tr}{\operatorname{Tr}}
\newcommand{\ad}{\operatorname{ad}}
\newcommand{\Vol}{\operatorname{Vol}}
\newcommand{\Spec}{\operatorname{Spec}}
\newcommand{\lam}{\lambda}
\newcommand{\om}{\omega}
\newcommand{\half}{\tfrac{1}{2}}
\newcommand{\sixth}{\tfrac{1}{6}}

\title{States of Low Energy on Bianchi II Spacetimes:\\
A Hadamard Construction via Nilpotent Plancherel Decomposition\\[6pt]
{\large \textit{Draft -- Gaps Not Closed}}}

\author{[Agent A5\_BII\_heisenberg, prepared for ECI v6.0.24+ project]}

\date{2026-05-03}

\begin{document}
\maketitle

\begin{abstract}
We initiate the construction of States of Low Energy (SLE) for a conformally
coupled massless scalar field on Bianchi type II cosmological backgrounds.
The spatial slices carry a left-invariant Riemannian metric on the Heisenberg group
$H_3$; after passage to a cocompact quotient, they become the compact nilmanifold
$\Nil = \Gamma \backslash H_3$.
We replace the Fourier decomposition used in the Bianchi I construction of
Banerjee--Niedermaier~\cite{BN23} by the Plancherel decomposition of
$L^2(H_3)$ via Stone--von Neumann theory, adapting the SLE energy functional
to integrate over the Plancherel measure $|\om|\,d\om$ on the dual $\widehat{H_3}$.
The spatial mode equation in each Plancherel sector $\om\in\ZZ\setminus\{0\}$
reduces to a time-dependent harmonic oscillator in the Schrödinger model coordinate,
with eigenvalues $\lam_{n,\om}(t) = (2n+1)|\om|/(a_1 a_2) + \om^2/a_3^2$.
The WKB adiabatic regime is controlled by the same large-$|\om|$ expansion as in the
flat case.
We identify the principal obstruction to a rigorous Hadamard theorem: a
uniformity-in-$n$ estimate for the Bogoliubov coefficients in the two-index
Plancherel sum, which does not follow from \cite{BN23} alone.
We conclude that the construction is algebraically complete but analytically open,
with an estimated time-to-publication of 6--9 months conditional on closing this gap.
\end{abstract}

\tableofcontents

\section{Introduction}

The construction of Hadamard states for quantum fields on cosmological backgrounds
with spatial homogeneity is a central problem in algebraic quantum field theory (AQFT)
in curved spacetime.
For Friedmann--Lemaître--Robertson--Walker (FLRW) spacetimes, adiabatic vacua and
their Hadamard property have been studied extensively;
for anisotropic Bianchi I backgrounds, the systematic SLE construction was carried out
by Banerjee--Niedermaier~\cite{BN23}, yielding the first rigorous Hadamard states
on the full Bianchi I family.

The natural next step in the Bianchi classification is Bianchi type II,
the simplest non-abelian unimodular case.
Its spatial symmetry group is the \emph{Heisenberg group} $H_3$, the prototype
of a non-flat nilpotent Lie group.
Compact spatial slices are the \emph{Heisenberg nilmanifold} (Nil-manifold)
$\Nil = \Gamma \backslash H_3$, the simplest compact 3-manifold beyond the
Bieberbach flat manifolds.

The goal of this note is to transfer the BN23 construction from the abelian
Bianchi I setting (where spatial symmetry is $\RR^3$ with flat Plancherel
$d^3 k$) to the nilpotent Bianchi II setting (spatial symmetry $H_3$,
Plancherel measure $|\om|\,d\om$ by Stone--von Neumann).
We carry out the algebraic and analytic setup in detail and identify the
remaining open problem for a complete Hadamard theorem.

\paragraph{Convention.}
Signature $(-,+,+,+)$, units $c = \hbar = 1$, conformal coupling $\xi = 1/6$.

\section{The Heisenberg Group and Its Representations}

\subsection{Lie algebra and group}

The \emph{Heisenberg Lie algebra} $\HH_3$ has basis $\{e_1, e_2, e_3\}$ with
\begin{equation}\label{eq:bracket}
  [e_1, e_2] = e_3, \qquad [e_1, e_3] = [e_2, e_3] = 0.
\end{equation}
The center is $\mathfrak{z} = \RR e_3$.
The adjoint representation matrices are (in the ordered basis $e_1, e_2, e_3$):
\begin{equation}
  \ad_{e_1} = \begin{pmatrix} 0&0&0\\ 0&0&0\\ 0&1&0 \end{pmatrix},
  \qquad
  \ad_{e_2} = \begin{pmatrix} 0&0&0\\ 0&0&0\\ -1&0&0 \end{pmatrix},
  \qquad
  \ad_{e_3} = 0.
\end{equation}

\begin{proposition}[Sympy-verified]\label{prop:unimod}
  $\Tr(\ad_X) = 0$ for all $X \in \HH_3$.
  In particular, $H_3$ is \emph{unimodular} (Bianchi class A).
  Moreover $(\ad_X)^2 = 0$ for all $X$, so $\HH_3$ is \emph{2-step nilpotent}.
\end{proposition}
\begin{proof}
  By direct matrix computation: $\Tr(\ad_{e_1}) = \Tr(\ad_{e_2}) = \Tr(\ad_{e_3}) = 0$.
  For general $X = a e_1 + b e_2 + c e_3$:
  $(a\,\ad_{e_1} + b\,\ad_{e_2})^2 = 0$ (verified by SymPy; see \texttt{sympy\_check.py}).
\end{proof}

The \emph{Heisenberg group} $H_3$ is the simply connected Lie group with Lie algebra $\HH_3$.
In coordinates $(x,y,z)$ given by $\exp(xe_1 + ye_2 + ze_3)$ (first-order Baker--Campbell--Hausdorff),
the group law is
\begin{equation}\label{eq:grouplaw}
  (x,y,z)\cdot(x',y',z') = (x+x',\; y+y',\; z+z'+xy').
\end{equation}

The \emph{left-invariant vector fields} (dual to the left Maurer--Cartan form) are
\begin{equation}\label{eq:livf}
  X_1 = \partial_x, \qquad
  X_2 = \partial_y + x\,\partial_z, \qquad
  X_3 = \partial_z,
\end{equation}
satisfying $[X_1, X_2] = X_3$ (confirms \eqref{eq:bracket}).
The dual left-invariant coframe is
\begin{equation}\label{eq:coframe}
  \sigma^1 = dx, \qquad \sigma^2 = dy, \qquad \sigma^3 = dz - x\,dy.
\end{equation}
One verifies $\sigma^i(X_j) = \delta^i_j$ directly.

\subsection{The Schrödinger representation and Plancherel theorem}

For each $\om \in \RR \setminus \{0\}$, define the \emph{Schrödinger representation}
$\piom: H_3 \to \mathcal{U}(L^2(\RR))$ by
\begin{equation}\label{eq:schrodinger_rep}
  (\piom(x,y,z)\psi)(q) = e^{i\om z}\, e^{i\om y q}\, \psi(q+x).
\end{equation}
This is unitary, irreducible, and has central character $\piom(\exp(ze_3)) = e^{i\om z}\cdot\mathrm{Id}$.

\begin{proposition}[Schrödinger rep is a Lie group homomorphism]
  The infinitesimal generators satisfy
  \begin{equation}
    d\piom(e_1) = \partial_q, \quad
    d\piom(e_2) = i\om q\cdot\mathrm{Id}, \quad
    d\piom(e_3) = i\om\cdot\mathrm{Id},
  \end{equation}
  and $[d\piom(e_1), d\piom(e_2)] = d\piom(e_3)$, consistent with $[e_1,e_2]=e_3$.
\end{proposition}
\begin{proof}
  $[d\piom(e_1), d\piom(e_2)]\psi(q) = [\partial_q, i\om q]\psi = i\om\psi = d\piom(e_3)\psi$,
  since $\partial_q(q\psi) - q\partial_q\psi = \psi$.
  (Verified by SymPy; see \texttt{sympy\_check.py}.)
\end{proof}

\begin{theorem}[Plancherel theorem for $H_3$,
  Stone--von Neumann~\cite{Folland1989}, Pukanszky~\cite{Pukan67}]\label{thm:plancherel}
  The unitary dual $\widehat{H_3}$ consists (up to measure zero) of the
  family $\{\piom : \om \in \RR\setminus\{0\}\}$ of Schrödinger representations,
  together with the trivial 1-dimensional representations $\pi_{p,q}$ ($p,q\in\RR$)
  of measure zero.
  The Plancherel decomposition is
  \begin{equation}\label{eq:plancherel}
    L^2(H_3) \cong \int_{\RR\setminus\{0\}} \HS(L^2(\RR))\, |\om|\,d\om,
  \end{equation}
  where $\HS(L^2(\RR))$ denotes Hilbert--Schmidt operators on $L^2(\RR)$.
  Explicitly, for $f \in L^1 \cap L^2(H_3)$:
  \begin{equation}
    \|f\|^2_{L^2} = \int_{\RR\setminus\{0\}} \|\piom(f)\|^2_{\HS}\,|\om|\,d\om,
    \qquad
    \piom(f) = \int_{H_3} f(h)\,\piom(h)\,d_L h.
  \end{equation}
\end{theorem}

\begin{remark}
  Since $H_3$ is unimodular (Proposition~\ref{prop:unimod}), the Plancherel measure
  $|\om|\,d\om$ is symmetric and there is no Duflo--Moore operator correction.
  The weight $|\om|$ arises from the Kirillov orbit method: the coadjoint orbits
  at $\om \neq 0$ are 2-dimensional planes in $\HH_3^*$, with symplectic area $|\om|$.
\end{remark}

\subsection{Cocompact quotient and quantization}

Fix the \emph{Heisenberg integer lattice}
$\Gamma = \{(m,n,p) \in H_3 : m,n,p \in \ZZ\}$.
One checks that $\Gamma$ is a cocompact lattice; the nilmanifold
$\Nil = \Gamma \backslash H_3$ has volume $\Vol(\Nil) = 1$.

On $\Nil$, the central character $\om$ is quantized:
the element $(0,0,1) \in \Gamma$ requires $e^{i\om \cdot 1} = 1$,
hence $\om \in \ZZ$.
The Plancherel integral~\eqref{eq:plancherel} restricts to a discrete sum,
\begin{equation}\label{eq:plancherel_nil}
  L^2(\Nil) \cong \bigoplus_{\om \in \ZZ\setminus\{0\}} V_\om,
  \qquad V_\om \cong L^2(\RR) \otimes \ell^2(\ZZ) \quad (\text{Hermite basis}),
\end{equation}
with the measure $|\om|\,d\om$ replaced by the weight $|\om|$ on each summand.

\section{Bianchi II Cosmological Background}

\subsection{Metric and curvature}

The Bianchi type II metric on $\RR \times \Nil$ is
\begin{equation}\label{eq:BII_metric}
  ds^2 = -dt^2 + a_1(t)^2\,(\sigma^1)^2 + a_2(t)^2\,(\sigma^2)^2 + a_3(t)^2\,(\sigma^3)^2,
\end{equation}
where $\sigma^i$ are the left-invariant 1-forms \eqref{eq:coframe} and
$a_i: \RR \to (0,\infty)$ are smooth scale factors.
In coordinates $(t,x,y,z)$ on $\RR \times H_3$, the metric components are
\begin{equation}
  g_{tt} = -1, \quad g_{xx} = a_1^2, \quad
  g_{yy} = a_2^2 + a_3^2 x^2, \quad
  g_{yz} = -a_3^2 x, \quad g_{zz} = a_3^2.
\end{equation}

\begin{lemma}[Sympy-verified]\label{lem:sqrt_g}
  $\det(g_{\mu\nu}) = -a_1^2 a_2^2 a_3^2$, hence
  $\sqrt{-g} = a_1(t)\,a_2(t)\,a_3(t)$,
  \emph{independent of the spatial coordinates $(x,y,z)$}.
\end{lemma}
\begin{proof}
  $\det(g_{ij}) = a_1^2\bigl[(a_2^2 + a_3^2 x^2)a_3^2 - a_3^4 x^2\bigr]
  = a_1^2 a_2^2 a_3^2$.
  (Verified by SymPy.)
\end{proof}

Lemma~\ref{lem:sqrt_g} is the analytic reflection of unimodularity:
the Haar measure on $H_3$ is the flat measure $dx\,dy\,dz$,
and the volume form on spatial slices is $a_1 a_2 a_3\,dx\,dy\,dz$.

The Ricci scalar for \eqref{eq:BII_metric} is (Kramer et al.\ 1980, standard Bianchi II curvature):
\begin{equation}\label{eq:ricci}
  R = -2\!\left(\frac{\ddot{a}_1}{a_1} + \frac{\ddot{a}_2}{a_2}
    + \frac{\ddot{a}_3}{a_3}\right)
    - 2\!\left(\frac{\dot{a}_1\dot{a}_2}{a_1 a_2}
    + \frac{\dot{a}_1\dot{a}_3}{a_1 a_3}
    + \frac{\dot{a}_2\dot{a}_3}{a_2 a_3}\right)
    + \frac{a_1^2}{2 a_2^2 a_3^2}.
\end{equation}
The last term $a_1^2/(2a_2^2 a_3^2)$ is the \emph{Heisenberg curvature contribution},
absent in the Bianchi I case ($a_1^2/(2a_2^2 a_3^2) = 0$ for flat spatial slices).

\subsection{Conformally coupled Klein--Gordon equation}

The conformally coupled massless scalar field satisfies
\begin{equation}\label{eq:KG}
  \left(\Box - \frac{1}{6}R\right)\phi = 0,
\end{equation}
where $\Box = (-g)^{-1/2}\partial_\mu\bigl((-g)^{1/2} g^{\mu\nu}\partial_\nu\bigr)$.
Using Lemma~\ref{lem:sqrt_g}, the wave operator splits as
\begin{equation}\label{eq:box_split}
  \Box\phi = -\ddot{\phi} - H_{\mathrm{eff}}\,\dot{\phi} + \Delta_t\phi,
\end{equation}
where $H_{\mathrm{eff}} = \dot{a}_1/a_1 + \dot{a}_2/a_2 + \dot{a}_3/a_3$
is the sum of Hubble rates and
\begin{equation}\label{eq:delta_t}
  \Delta_t = a_1^{-2}X_1^2 + a_2^{-2}X_2^2 + a_3^{-2}X_3^2
\end{equation}
is the time-dependent spatial Laplacian built from the left-invariant vector fields
$X_i$ in \eqref{eq:livf}.

\section{Plancherel Mode Decomposition}

\subsection{Harmonic analysis in the $\pi_\omega$ sector}

Fix $\om \in \ZZ \setminus \{0\}$ and consider the $\om$-isotypic component of the field.
The central element $X_3 = \partial_z$ acts on this component as $i\om$.

After a partial Fourier transform in $y$ (with momentum $k_y \in \RR$),
the spatial Laplacian~\eqref{eq:delta_t} restricted to the $(\om, k_y)$ sector becomes
\begin{equation}\label{eq:laplacian_sector}
  \Delta_t^{(\om, k_y)} = a_1^{-2}\partial_x^2
    - a_2^{-2}(k_y + \om x)^2
    - a_3^{-2}\om^2.
\end{equation}
This is a \emph{1D quantum harmonic oscillator} in $x$,
centered at $x_0 = -k_y/\om$,
with time-dependent angular frequency
\begin{equation}\label{eq:osc_freq}
  \Omega(t) = \frac{|\om|}{a_1(t)\,a_2(t)}.
\end{equation}

\begin{proposition}\label{prop:eigenvalues}
  The eigenfunctions of $-\Delta_t^{(\om,k_y)}$ in $L^2(\RR_x)$ are
  the shifted Hermite functions $\psi_{n}^{(\om)}(x) = h_n(\sqrt{\Omega}\,(x-x_0))$
  with eigenvalues
  \begin{equation}\label{eq:eigenvalues}
    \lam_{n,\om}(t) = (2n+1)\,\frac{|\om|}{a_1(t)\,a_2(t)} + \frac{\om^2}{a_3(t)^2},
    \qquad n \in \NN_0.
  \end{equation}
\end{proposition}
\begin{proof}
  Completing the square: $-a_2^{-2}(k_y+\om x)^2 = -\om^2 a_2^{-2}(x-x_0)^2$.
  The operator $-a_1^{-2}\partial_x^2 + \om^2 a_2^{-2}(x-x_0)^2$ is a harmonic
  oscillator with frequency $\Omega = |\om|/(a_1 a_2)$, eigenvalues $(2n+1)\Omega$.
  Adding $-a_3^{-2}\om^2$ (constant in $x$) gives~\eqref{eq:eigenvalues}.
\end{proof}

\paragraph{UV behavior.}
For $|\om| \to \infty$:
\begin{equation}\label{eq:UV_behavior}
  \lam_{n,\om}(t) = \frac{\om^2}{a_3^2} + (2n+1)\frac{|\om|}{a_1 a_2} + O(1).
\end{equation}
The \emph{leading} UV behavior is $\om^2/a_3^2$, identical to the flat (Bianchi I)
spectrum $k_3^2/a_3^2$ with $k_3 \leftrightarrow \om$.
The subleading term $(2n+1)|\om|/(a_1 a_2)$ is the Heisenberg non-commutativity contribution.

\subsection{Mode expansion}

The field is expanded in the Plancherel--Hermite basis:
\begin{equation}\label{eq:mode_expansion}
  \phi(t,x,y,z) = \sum_{\om \in \ZZ\setminus\{0\}} |\om|
    \int_\RR \sum_{n=0}^\infty
    \phi_{n,\om}^{(k_y)}(t)\,
    e^{i k_y y}\,e^{i\om z}\,
    h_n\!\left(\frac{|\om|^{1/2}}{a_1^{1/2} a_2^{1/2}}\bigl(x + \tfrac{k_y}{\om}\bigr)\right)
    \frac{dk_y}{2\pi},
\end{equation}
where $h_n$ are the normalized Hermite functions.

Each complex amplitude $\phi_{n,\om}^{(k_y)}(t)$ satisfies the \emph{mode equation}
\begin{equation}\label{eq:mode_eq}
  \ddot{\phi}_{n,\om} + H_{\mathrm{eff}}\,\dot{\phi}_{n,\om}
  + \omega_{n,\om}^2(t)\,\phi_{n,\om} = 0,
\end{equation}
with total squared frequency
\begin{equation}\label{eq:total_freq}
  \omega_{n,\om}^2(t) = \lam_{n,\om}(t) + \sixth R(t)
  = (2n+1)\frac{|\om|}{a_1 a_2} + \frac{\om^2}{a_3^2} + \sixth R(t).
\end{equation}

\paragraph{Conformal redefinition.}
Setting $v_{n,\om} = (a_1 a_2 a_3)^{1/2}\,\phi_{n,\om}$ and using conformal time,
the mode equation~\eqref{eq:mode_eq} reduces to a Schrödinger-type equation
\begin{equation}\label{eq:schrodinger_mode}
  v_{n,\om}'' + \bigl(\omega_{n,\om}^2(\eta) - m_{\mathrm{eff}}^2(\eta)\bigr)\,v_{n,\om} = 0,
\end{equation}
where primes denote conformal-time derivatives and
$m_{\mathrm{eff}}^2 = (a_1 a_2 a_3)^{1/2}\,[(a_1 a_2 a_3)^{-1/2}]''$
is the standard effective mass correction.

\section{The SLE Functional on Bianchi II}

\subsection{Definition}

Following Banerjee--Niedermaier~\cite{BN23}, fix a real-valued test function
$f \in C_c^\infty(\RR_t)$ (temporal smearing).
Let $T_{00}$ be the energy density of the conformally coupled field~\eqref{eq:KG}.
Define the smeared energy functional
\begin{equation}\label{eq:SLE_functional}
  \mathcal{F}_f[\phi] = \int_\RR |f(t)|^2
    \int_{\Nil} T_{00}(t,h)\,\sqrt{\gamma}\,dh\,dt,
\end{equation}
where $\sqrt{\gamma} = a_1 a_2 a_3$ is the spatial volume element.
After inserting the mode expansion~\eqref{eq:mode_expansion}:

\begin{equation}\label{eq:SLE_functional_modes}
  \mathcal{F}_f[\phi] = \sum_{\om \in \ZZ\setminus\{0\}} |\om|
    \int_\RR \sum_{n=0}^\infty
    \int_\RR |f(t)|^2\,
    \Bigl[|\dot{\phi}_{n,\om}|^2 + \omega_{n,\om}^2(t)\,|\phi_{n,\om}|^2\Bigr]
    dt\;\frac{dk_y}{2\pi}.
\end{equation}

\begin{proposition}[Well-definedness]\label{prop:well_def}
  For $f \in C_c^\infty(\RR)$ and $\phi$ in the classical solution space,
  the functional~\eqref{eq:SLE_functional_modes} is well-defined provided
  the initial data for $\phi_{n,\om}$ satisfy the WKB normalization
  $|\phi_{n,\om}|^2 \sim (2\omega_{n,\om})^{-1}$ as $|\om|,n \to \infty$.
\end{proposition}
\begin{proof}[Proof sketch]
  By \eqref{eq:UV_behavior}, $\omega_{n,\om}^2 \sim \om^2/a_3^2$ for large $|\om|$.
  The contribution of each mode to $\mathcal{F}_f$ is proportional to
  $|\om| \cdot \omega_{n,\om} \sim |\om|^2$.
  Summing with WKB weight $\omega_{n,\om}^{-1} \sim a_3/|\om|$ and Plancherel weight $|\om|$:
  the sum over $n$ converges absolutely for fixed $\om$ (the harmonic oscillator
  ground-state energy gives $(2n+1)\Omega \cdot (2(2n+1)\Omega)^{-1} = 1/2$,
  but the Hermite coefficients of $f$ provide the damping needed for the $n$-sum).
  Convergence in $\om$ follows from the rapid decay of Plancherel--Fourier coefficients
  of smooth functions on $\Nil$.
  A rigorous proof requires the uniformity-in-$n$ estimate discussed in \S\ref{sec:gaps}.
\end{proof}

\subsection{Minimization and the SLE state}

The mode equation~\eqref{eq:mode_eq} has a 2-dimensional solution space for each $(\om,n,k_y)$.
The SLE state is defined by choosing, for each mode, the solution that
minimizes $\mathcal{F}_f$, subject to the Wronskian normalization
$\phi_{n,\om}\dot{\bar{\phi}}_{n,\om} - \dot{\phi}_{n,\om}\bar{\phi}_{n,\om} = -i$.

\begin{proposition}[Adaptation of BN23 \S3]\label{prop:SLE_mode}
  For each sector $(\om,n)$, the minimizer of $\int |f(t)|^2
  [|\dot\phi|^2 + \omega_{n,\om}^2|\phi|^2]dt$ subject to the Wronskian constraint
  satisfies the Euler--Lagrange equation
  \begin{equation}\label{eq:EL}
    \ddot{u} + \omega_{n,\om}^2(t)\,u = 0
  \end{equation}
  with initial conditions at a reference time $t_0$ given by the WKB approximation:
  \begin{equation}\label{eq:WKB_IC}
    u(t_0) = \frac{1}{\sqrt{2\omega_{n,\om}(t_0)}},
    \qquad
    \dot{u}(t_0) = -i\sqrt{\frac{\omega_{n,\om}(t_0)}{2}}.
  \end{equation}
\end{proposition}
\begin{proof}
  Identical to BN23~\cite{BN23} Proposition~3.2, replacing $k \mapsto (\om,n)$
  and $\omega_k^2 \mapsto \omega_{n,\om}^2(t)$.
  The mode equation~\eqref{eq:mode_eq} is formally identical after the
  conformal redefinition.
\end{proof}

\section{Towards the Hadamard Property}\label{sec:hadamard}

\subsection{Radzikowski's criterion}

By Radzikowski~\cite{Rad96}, a quasi-free state $\omega$ for a Klein--Gordon field
on a globally hyperbolic spacetime has the Hadamard property if and only if
\begin{equation}\label{eq:WF_Hadamard}
  \WF(\omega_2) = \{(x,k;x',k') : (x,k) \sim_+ (x',k'),\; k \text{ future-directed null}\},
\end{equation}
where $\omega_2$ is the two-point distribution and $\sim_+$ denotes the relation
of being connected by a null geodesic with positive-frequency propagation.

\subsection{Two-point function in the SLE state}

In the SLE state defined by~\eqref{eq:WKB_IC}, the two-point function is
\begin{equation}\label{eq:2pt_function}
  \omega_2(t,h;\,t',h') = \sum_{\om\in\ZZ\setminus\{0\}} |\om|
    \sum_{n=0}^\infty \int_\RR
    u_{n,\om}(t)\,\overline{u_{n,\om}(t')}\,
    \mathcal{K}_{n,\om}^{(k_y)}(h,h')\,\frac{dk_y}{2\pi},
\end{equation}
where $u_{n,\om}$ solves~\eqref{eq:EL}--\eqref{eq:WKB_IC} and
$\mathcal{K}_{n,\om}^{(k_y)}(h,h')$ is the reproducing kernel for the
$(n,\om,k_y)$-eigenfunction.

\subsection{Expected Hadamard argument (partial)}

A complete Hadamard proof would proceed in three steps:

\textbf{Step H1 (temporal):}
For fixed $(\om,n)$, the function $u_{n,\om}(t)\,\overline{u_{n,\om}(t')}$
is the two-point function of a Hadamard state for a 1D harmonic oscillator
with time-dependent frequency $\omega_{n,\om}(t)$.
By the adiabatic theorem and BN23~\cite{BN23} Theorem~3.1,
$\WF(u_{n,\om}\,\bar{u}_{n,\om})$ lies in the null bicharacteristic relation
of the 1+1-dimensional operator $-\partial_t^2 + \omega_{n,\om}^2(t)$.

\textbf{Step H2 (spatial):}
The spatial kernel $\mathcal{K}_{n,\om}^{(k_y)}(h,h')$ is smooth off the diagonal
and has a Riemannian Hadamard singularity on $h=h'$,
controlled by the heat kernel on $(\Nil, g_t)$ where $g_t$ is the spatial metric
at fixed $t$.
The Riemannian heat kernel on Nil$^3$ is known; see the eigenvalue formula
in Proposition~\ref{prop:eigenvalues}.

\textbf{Step H3 (uniformity):}
The sum~\eqref{eq:2pt_function} over $(\om, n, k_y)$ converges in the distributional
sense with WF set given by~\eqref{eq:WF_Hadamard}.
This requires \emph{uniform decay of Bogoliubov coefficients}:
\begin{equation}\label{eq:beta_decay}
  |\beta_{n,\om}|^2 \leq C_N\,(|\om| + n)^{-N} \qquad \forall N \in \NN,
\end{equation}
uniformly in both indices $\om$ and $n$.

Step H3 is the \emph{principal open problem} and is discussed in \S\ref{sec:gaps}.

\section{Obstructions and Open Problems}\label{sec:gaps}

\subsection{Gap 1: Wavefront set uniformity (hard, blocks the theorem)}

Estimate~\eqref{eq:beta_decay} is the two-index generalization of the Bogoliubov
coefficient decay used in BN23~\cite{BN23} Theorem~3.1 and Appendix~A.
For Bianchi I (abelian, single index $k \in \ZZ^3$), the adiabatic expansion
gives $|\beta_k|^2 = O(|k|^{-N})$ for all $N$ by smooth dependence of the
frequency $\omega_k(t)$ on $t$.

For Bianchi II, the two-index structure introduces a new subtlety:
for fixed $\om$, the frequency $\omega_{n,\om}(t) = [(2n+1)\Omega(t) + \om^2/a_3^2 + \sixth R]^{1/2}$
grows as $n^{1/2}$ for large $n$ (fixed $\om$).
The adiabatic expansion for large $n$ must be controlled \emph{uniformly in $\om$}.
This is a distinct problem from the large-$|\om|$ adiabaticity and requires
a two-dimensional asymptotic analysis.
Standard references (Avron--Seiler--Yaffe 1987~\cite{ASY87}, Reed--Simon Vol.~IV~\cite{RS4})
give the result for fixed $\om$, but the uniformity in $\om$ is not standard.

\textbf{Resolution path:}
Adapt the parametrix construction of Junker--Schrohe~\cite{JS02}
(adiabatic vacua of higher order) to the two-index setting,
treating $(\om, n)$ as a joint semiclassical parameter.

\subsection{Gap 2: Nil$^3$ Weyl law and Sobolev threshold}

The spatial Weyl counting function on $\Nil = \Gamma \backslash H_3$
satisfies $N(\lam) \sim C\lam^2$ (not $\lam^{3/2}$ as for flat $T^3$),
because the dominant growth of eigenvalues~\eqref{eq:eigenvalues} is
$\lam_{n,\om} \sim \om^2$ for large $|\om|$ at fixed $n$.
This changes the Sobolev embedding threshold.
The SLE functional~\eqref{eq:SLE_functional_modes} converges for
smearing functions $f \in H^2(\Nil)$ rather than $H^1$ as in the flat case.

\subsection{Gap 3: Absent timelike Killing vector}

On Bianchi II, there is no global timelike Killing vector.
The SLE state depends on the choice of smearing function $f$ and reference time $t_0$.
Stability of the Hadamard property under change of $t_0$ follows formally from the
adiabatic expansion but requires the uniformity of Gap~1.

\subsection{Gap 4: IR regularity ($\om \to 0$)}

On the nilmanifold, $\om \in \ZZ\setminus\{0\}$ so $|\om| \geq 1$ and
the IR is automatically regulated.
On the universal cover $H_3$ (non-compact spatial sections),
the $\om \to 0$ limit requires separate analysis of the abelian sector.
For compact $\Nil$, this is not an issue.

\subsection{Gap 5: Connection to ECI framework}

The Bianchi II SLE state should be mapped into the ECI lever formalism.
The $\om$-index (Plancherel central character, physically the ``magnetic quantum number''
of the Heisenberg orbit) may correspond to a new ECI observable.
This is entirely open.

\section{Conclusions}

We have established:
\begin{enumerate}[label=(\roman*)]
\item The Heisenberg group $H_3$ is unimodular with Plancherel measure $|\om|\,d\om$
  (Proposition~\ref{prop:unimod}, Theorem~\ref{thm:plancherel}).
\item The spatial mode decomposition on Bianchi II reduces each Plancherel sector
  $(\om, n)$ to a time-dependent harmonic oscillator with eigenvalues~\eqref{eq:eigenvalues}.
\item The SLE energy functional~\eqref{eq:SLE_functional_modes} is formally well-defined
  and the minimizer satisfies mode equations~\eqref{eq:EL}--\eqref{eq:WKB_IC} identical
  in form to BN23.
\item The UV behavior $\lam_{n,\om} \sim \om^2/a_3^2$ is the same as in the flat case,
  ensuring the WKB regime is controlled.
\item The Hadamard property reduces to a wavefront set uniformity estimate (Gap~1)
  that is plausible but not yet proved.
\end{enumerate}

\textbf{Verdict:} The construction is algebraically complete and structurally sound.
It does \emph{not yet} constitute a rigorous Hadamard theorem.
The piste remains ``Type A: rigorous pending Hadamard,'' with an estimated
time-to-publication of \textbf{6--9 months} (see \texttt{gaps.md}).

\bigskip
\noindent\textbf{Note on one hallucinated reference:}
The task brief cited ``Pukanszky 1967 TAMS 126, 487-507.''
This reference does not exist.
The correct citation is J.\ Funct.\ Anal.\ \textbf{1} (1967) 255--280~\cite{Pukan67}.

\section*{Acknowledgments}
Prepared as output of agent A5\_BII\_heisenberg for the ECI v6.0.24+
piste on Bianchi II cosmological states.

\begin{thebibliography}{99}

\bibitem{BN23}
R.~Banerjee, M.~Niedermaier,
\textit{States of Low Energy on Bianchi I spacetimes},
J.\ Math.\ Phys.\ \textbf{64} (2023) 113503.
arXiv:2305.11388.

\bibitem{AV13}
Z.~Avetisyan, R.~Verch,
\textit{Explicit harmonic and spectral analysis in Bianchi I-VII type cosmologies},
Class.\ Quant.\ Grav.\ \textbf{30} (2013) 155006.
arXiv:1212.6180.

\bibitem{Pukan67}
L.~Pukánszky,
\textit{On the characters and the Plancherel formula of nilpotent groups},
J.\ Funct.\ Anal.\ \textbf{1} (1967) 255--280.
DOI:10.1016/0022-1236(67)90015-8.
% VERIFIED. NOT Trans. Amer. Math. Soc. (that reference is wrong/hallucinated).

\bibitem{Folland1989}
G.~B.~Folland,
\textit{Harmonic Analysis in Phase Space},
Annals of Mathematics Studies \textbf{122}, Princeton University Press, 1989.
ISBN~978-0-691-08528-9.
% VERIFIED: publisher, series, year confirmed.

\bibitem{Rad96}
M.~J.~Radzikowski,
\textit{Micro-local approach to the Hadamard condition in quantum field theory
on curved space-time},
Commun.\ Math.\ Phys.\ \textbf{179} (1996) 529--553.
DOI:10.1007/BF02100096.
% VERIFIED via Springer link.

\bibitem{SvN}
J.~von Neumann,
\textit{Die Eindeutigkeit der Schrödingerschen Operatoren},
Math.\ Ann.\ \textbf{104} (1931) 570--578.
% Standard Stone-von Neumann theorem; classical result.

\bibitem{JS02}
W.~Junker, E.~Schrohe,
\textit{Adiabatic vacuum states on general spacetime manifolds:
definition, construction, and physical properties},
Ann.\ Henri Poincaré \textbf{3} (2002) 1113--1181.
% Adiabatic vacua, relevant for Bogoliubov coefficient decay.
% NOTE: To be verified against arXiv before final submission.

\bibitem{ASY87}
J.~E.~Avron, R.~Seiler, L.~G.~Yaffe,
\textit{Adiabatic theorems and applications to the quantum Hall effect},
Commun.\ Math.\ Phys.\ \textbf{110} (1987) 33--49.
% Adiabatic theorem with gap condition; used for WKB uniformity argument.
% NOTE: To be verified against DOI before final submission.

\bibitem{RS4}
M.~Reed, B.~Simon,
\textit{Methods of Modern Mathematical Physics, Vol.~IV: Analysis of Operators},
Academic Press, 1978.
% Standard reference for adiabatic theorem in spectral theory.

\bibitem{Mal51}
A.~I.~Mal'cev,
\textit{On a class of homogeneous spaces},
Amer.\ Math.\ Soc.\ Transl.\ \textbf{39} (1951).
% Lattice existence in nilpotent Lie groups.
% NOTE: Verify translation reference before submission.

\end{thebibliography}

\end{document}
